\documentclass[11pt]{amsart}

\usepackage{lmodern}
\usepackage{microtype}
\usepackage[english]{babel}
\usepackage{mathtools}
\usepackage{amssymb}
\usepackage{booktabs}
\usepackage[a4paper,margin=28mm]{geometry}
\usepackage[hidelinks]{hyperref}
\usepackage[msc-links]{amsrefs}
\usepackage[capitalize,nameinlink,noabbrev]{cleveref}

\begin{document}

\title[Mapping a research area with AI agents]{A map of the range, not a way
  up the wall: six weeks of AI agents on 373 open problems in analysis}
\author{Øistein Søvik}
\date{\today}

\begin{abstract}
For six weeks, one mathematician with a master's degree directed AI coding
agents at every open problem that could be found in one area of analysis:
Hardy, Bergman and Dirichlet-series spaces, Hankel and Helson operators, and
Nehari-type theorems. The agents collected 373 problems, found that 64 of them
had already been settled by others, and produced answers to 141 more. We
classify these answers by how they were obtained and by how much they matter.
Most are easy in hindsight: short arguments from known tools, results assembled
from the literature, or counterexamples that are easy to check once found. Six
are, in our judgement, of real significance. On the hardest problems, days of
work by several agents produced maps of what fails, not proofs. We conclude
that current agents are very good at mapping a field and clearing its long
tail, and that a hard problem stays hard. All results of ours are sketches
until they appear on arXiv or in a journal; this is a living document, updated
as the people who asked the questions respond.
\end{abstract}

\maketitle

\section{Introduction}

Terence Tao has described research mathematics as a range of mountains:
\begin{quote}
  Tao gave the analogy of mathematicians trying to climb ``a big mountain
  range with lots of tall mountains and lots of foothills''. Humans can only
  climb one step at a time, but they can plan a route to the top of a
  mountain like Everest. Meanwhile, Tao said, current AIs are like jumping
  robots. They can sometimes ``parkour their way to the top of a 6-foot
  wall'' that a human couldn't climb. But they can't do long-term strategic
  planning. Those 6 feet might become 10 feet, or 100, Tao imagines, but
  ``the little jumping robots are nowhere near the Mount Everests of
  math''~\cite{quanta}.
\end{quote}
This paper is a field report on what such robots do when pointed at an entire
range. Over six weeks in autumn 2026 we used Claude Code and Codex agents to
find, sort and attack every open problem we could locate in one research area.
The agents found 373 problems in papers, problem lists and monographs. They
established that 64 were already solved, and they produced answers to 141
others.

The answers need careful reading. Our main finding is that the gain was in
cartography, not in climbing:
\begin{itemize}
  \item The agents were very good at mapping a field: finding problems,
        classifying them, checking the literature, and clearing problems that
        were easy but had received little attention.
  \item Most of their answers are low-hanging fruit. Of the 141, we classify
        6 as significant, 62 as useful, and 73 as minor.
  \item On problems where experts had made partial progress and stopped,
        the agents produced reformulations, special cases and catalogues of
        obstructions, not proofs.
\end{itemize}
Breadth explains the counts. Throw a hundred balls at a hundred baskets and
some go in; that is useful, and different from getting better at the shot that
matters. Tao's description of current tools as ``artificial general
cleverness''~\cite{tao-cleverness} fits what we saw.

We therefore present this method as the first step of a research project, not
the last: a fast way to get to know a field and to find questions worth working
on. \Cref{sec:map} describes how the problems were found and sorted,
\cref{sec:workflow} the workflow, \cref{sec:results} the results and their
classification, \cref{sec:walls} the hard problems, and \cref{sec:verification}
verification and credit. \Cref{sec:related} places the work among other
reports, and \cref{sec:limits} lists its limits.

\section{Setting}

The area is the theory of Hardy and Bergman spaces on the disc and on the
infinite polydisc, Dirichlet series viewed as functions on the infinite torus,
and the Hankel, Helson and composition operators acting on them. A typical
question asks for the sharp constant in an inequality, whether an extremal
function can have zeros, or whether a bounded bilinear form has a bounded
symbol. The active literature is a few hundred papers, much of it from a small
number of research groups.

One person directed the work. He chose targets, judged whether a result
mattered, steered agents that were stuck, read everything before it went out,
and handled all contact with other mathematicians. The agents were Claude Code
(Claude Opus 5.5 as the main session, with smaller Claude models for cheap
subtasks) and Codex (several GPT models), often running at the same time in the
same repository. The record of the work is the version history, with about 1000
commits, together with the logs: 14 Claude Code sessions with 393 subagents,
and 586 Codex sessions, 420 of them subagents.

\section{Building the map}\label{sec:map}

The problems came from four sources, in this order.
\begin{enumerate}
  \item \emph{Stated problems in a local paper collection.} Agents searched
        each paper for environments labelled Conjecture, Question, Problem or
        Open problem, and for phrases such as ``remains open'' and ``it would be
        interesting to''. Each candidate was recorded with the paper, the exact
        location, and a quote.
  \item \emph{The papers those cite}, with the same extraction.
  \item \emph{Synthesis passes.} Agents read the recent papers of the core
        authors ``between the lines''. This produced \emph{implicit} questions,
        which a remark suggests without posing, and \emph{inferences}: paper A
        proves one statement and paper B another, so a third should follow.
  \item \emph{Problem lists and monographs}: workshop problem sessions, a
        classical problem book, and the open-problem sections of two monographs.
\end{enumerate}

Each problem became one Markdown file with validated metadata. The metadata
records:
\begin{itemize}
  \item the theme, the kind (298 stated, 55 implicit, 20 inferences) and the
        origin (289 asked by others, 84 by us);
  \item tags, an impact score and a difficulty score;
  \item a dated literature check with quoted evidence;
  \item the stage reached in our own work;
  \item once anyone settled the problem, how far the resolution is trusted:
        journal, arXiv, or our own.
\end{itemize}
Moving from a single catalogue file to one file per problem let several agents
and two different tools edit the database at once.

\section{The workflow}\label{sec:workflow}

Each problem was worked in sessions. A session begins by reading the record,
the source passage and the log of earlier attempts, and by writing a three-line
plan. It ends by rewriting a short \emph{state of play} and appending one line
to the \emph{attempt log}: the route, the outcome, the reason, and the next
step. A claimed result had to become a \emph{solution sheet}: the problem as
posed, the answer as a theorem, a complete proof with exact citations, and
nothing else, built through a lint-gated LaTeX pipeline. A fresh agent with a
fixed brief then refereed it. The referee re-read the question in its source,
re-derived each step, reran any certificate, and decided whether the answer was
full or partial.

On hard problems several agents worked in parallel \emph{lanes}, each on one
approach, with a coordinating session redirecting them. When lanes drifted into
cataloguing obstructions, a \emph{prove-mode} brief told each to commit to one
strategy and push it to a theorem, or to a proof complete except for one
precisely stated lemma. The full instructions are published with this paper as
a reusable set of agent playbooks.

\section{Results}\label{sec:results}

\subsection{What the literature check found}

The difficulty assigned to a problem before anyone attempted it predicted the
outcome well (\cref{tab:difficulty}). It also showed why easy problems are
easy: nine of the twelve problems first rated easiest already had a published
answer. Tao observed the same correlation on the Erdős problems: the problems
current tools can solve are ``precisely the category of nominally open problems
that are most likely to have been solved in the
literature''~\cite{tao-lowhanging}. The literature check also caught one
problem recorded as resolved by a paper that had since been withdrawn.

\begin{table}[ht]
  \centering
  \begin{tabular}{lrrrr}
    \toprule
    Initial difficulty & Settled by others & Settled by us & Partial & Open \\
    \midrule
    1                  & 9                 & 3             & 0       & 0    \\
    2                  & 4                 & 15            & 2       & 0    \\
    3                  & 7                 & 27            & 19      & 7    \\
    4                  & 4                 & 29            & 75      & 6    \\
    5                  & 0                 & 3             & 15      & 3    \\
    \bottomrule
  \end{tabular}
  \caption{Outcome by the difficulty assigned before any attempt, for the
    275 records that had one.}\label{tab:difficulty}
\end{table}

\subsection{What kind of results were these?}

We classified each of the 141 answers along three axes.
\begin{itemize}
  \item \emph{Sign}: 82 proofs, 42 counterexamples, and 17 determinations
        of a value, set or order.
  \item \emph{Kind}, meaning how the answer was obtained:
        \begin{itemize}
          \item 20 assembled from the literature;
          \item 58 routine new arguments;
          \item 44 explicit constructions or certificates, which are the work
                to find and easy to check;
          \item 19 substantial new arguments.
        \end{itemize}
  \item \emph{Significance}:
        \begin{itemize}
          \item 6 high: paper-level answers to someone else's question;
          \item 62 medium: full answers worth a short note;
          \item 73 low: our own questions, conditional or elementary answers,
                and corollaries.
        \end{itemize}
\end{itemize}
The classification is ours and preliminary; \cref{tab:sign} gives the
cross-tabulations.

\begin{table}[ht]
  \centering
  \begin{tabular}{lrrrrr}
    \toprule
                   & Literature & Routine & Construction & Substantial & Total \\
    \midrule
    Proof          & 17         & 43      & 5            & 17          & 82    \\
    Counterexample & 0          & 1       & 39           & 2           & 42    \\
    Determination  & 3          & 14      & 0            & 0           & 17    \\
    \bottomrule
  \end{tabular}

  \medskip
  \begin{tabular}{lrrrr}
    \toprule
                    & High & Medium & Low & Total \\
    \midrule
    Proof           & 4    & 33     & 45  & 82    \\
    Counterexample  & 2    & 24     & 16  & 42    \\
    Determination   & 0    & 5      & 12  & 17    \\
    \midrule
    Asked by others & 6    & 57     & 26  & 89    \\
    Asked by us     & 0    & 5      & 47  & 52    \\
    \bottomrule
  \end{tabular}
  \caption{The 141 answers by sign and kind (top) and by sign, origin and
    significance (bottom).}\label{tab:sign}
\end{table}

Three things stand out.
\begin{itemize}
  \item \emph{Counterexamples are the natural strength.} Thirty-nine of the
        42 counterexamples are explicit constructions or certificates. Searching
        a space of candidates and checking each one is exactly what agents do
        well, and what earlier work on machine-assisted discovery exploited
        \cites{wagner,funsearch,alphaevolve}.
  \item \emph{Few proofs are substantial.} Only 17 of the 82 proofs needed
        a genuinely new argument; the rest are routine or assembled from the
        literature.
  \item \emph{The significant results needed direction.} Of the six results
        we rate high, four are proofs and two are counterexamples. Three of the
        six came from focused sessions in which the human directed one problem
        with the asker's question in view, not from the bulk campaign over the
        database.
\end{itemize}

The six high-significance results are:
\begin{itemize}
  \item that extremal functions for bounded linear functionals on $H^p$,
        $0<p<1$, are outer, answering a published conjecture in a stronger
        form, written jointly with the asker;
  \item a sharp bound on the number of zeros of extremals in a
        prescribed-coefficient problem;
  \item the finite-exponent limits of sharp Bernstein--Nikolskii constants;
  \item the exact order of growth of the $H^1$ norms of partial-sum
        operators for Dirichlet series;
  \item a counterexample to a positive-symbol Nehari question, with a
        computer certificate in outward-rounded arithmetic;
  \item that optimal polynomial approximants in $H^p$ can vanish in the
        disc for every $p\neq2$.
\end{itemize}
Each has been sent to the people who asked the question, or is being prepared
for them. None is yet on arXiv.

\section{Walls}\label{sec:walls}

We attacked several problems on which experts have made partial progress, for
days and with several agents. Two examples, described without the technical
details:
\begin{itemize}
  \item \emph{An endpoint Schatten-class Nehari problem.} Every route the
        agents found reduced to one square-function estimate in $L^1$ that is a
        theorem in its own right. They proved many special cases, wrote down
        equivalent reformulations, and catalogued exactly where each route
        failed. They did not find a proof, and the problem was handed back to
        the person who asked it.
  \item \emph{A contractive restriction conjecture.} The agents proved the
        conjecture for two natural families of functions and refuted a stronger
        statement that would have implied it. The general case did not move.
\end{itemize}
In both cases the output was useful as a map. It is not progress on the wall
itself, and we do not present it as such. The pattern matches Tao's image: some
walls of six feet were climbed; the Everests were not approached.

\section{Verification and credit}\label{sec:verification}

Producing a candidate proof was cheap; checking it was the bottleneck. This
matches Tao's analysis of the problem-solving pipeline, in which generation is
followed by verification, exposition, acceptance and absorption into the
standard theory, and in which AI threatens ``proof
indigestion''~\cite{tao-age}. Our answers were checked by agent referees, in
some cases by a second model family, and in a few cases by Lean formalization
or exact-arithmetic certificates. Agent referees caught real errors: a bound
used outside its range, a missing hypothesis, and above all overstated scope,
where an answer to a special case was reported as an answer to the question.
They rarely asked whether the claim was the right one.

For this reason every result of ours is labelled a \emph{sketch} until it
appears on arXiv or in a journal, however it was checked. Before anything is
posted, the people who asked the question are told; this follows the
recommendations of the Leiden Declaration on attribution, disclosure and the
humanity of authorship~\cite{leiden}. The questions belong to those who asked
them. The proofs and computations were produced by AI agents and checked by the
author. Following Tao's rule of thumb~\cite{tao-age}, no result should be
published until a human can give a clear, expert-level talk on it.

\section{Related work}\label{sec:related}

\emph{Machine-assisted discovery.} Machine learning has guided human intuition
in knot theory and representation theory~\cite{davies}, and reinforcement
learning has found counterexamples to conjectures in
combinatorics~\cite{wagner}. Language models coupled with evaluators have found
new constructions~\cites{funsearch,alphaevolve}, and have been applied to many
problems at once~\cite{georgiev}. Our counterexamples are of this kind:
constructions that are hard to find and easy to check.

\emph{Formal proof.} Language models for theorem proving date back at least
to~\cites{polu,htps}; olympiad geometry has been solved without human
demonstrations~\cite{alphageometry}. The mathlib library~\cite{mathlib}
underlies much of this, and Avigad discusses the wider ``formal
turn''~\cite{avigad}.

\emph{Evaluation.} Early evaluations of language models on mathematics
include~\cites{frieder,collins}. Benchmarks with unpublished
problems~\cites{frontiermath,hle} and the First Proof project, which has
experts referee AI solutions to unpublished research
problems~\cite{firstproof}, measure capability under controlled conditions. Our
report is neither controlled nor a benchmark; it is one practitioner's account
of a whole area.

\emph{Research practice.} Experiments with frontier models as research
assistants are collected in~\cite{bubeck}. The Erdős problems
database~\cite{erdos} and its AI wiki~\cite{erdos-wiki} are the closest
analogue to our problem database; the wiki warns that ``absence of past
progress may reflect obscurity rather than difficulty''. On the goals of
mathematics beyond problem solving we follow Thurston~\cite{thurston},
Gowers~\cite{gowers}, Tao~\cites{tao-good,tao-age} and Bessis~\cite{bessis}.
Journalism on the current state of the field includes~\cite{quanta}, where Igor
Pak notes that language models are ``fantastic in finding the right
references''.

\section{Limits}\label{sec:limits}

This is one person's experience in one area, over six weeks, without a control
group. The difficulty and significance scores were assigned by the same agents
and by the author, not by independent experts. Correctness has not been audited
on a random sample; we intend to do so, and the counts above should be read as
counts of claimed answers. The models will be superseded within months.
Finally, the area has a dense literature of recent preprints, which made the
literature check unusually productive.

\section*{Tool and computational resource disclosure}

The work used Claude Code from Anthropic, with Claude Opus 5.5, Sonnet and
Haiku, and Codex from OpenAI, with several GPT models; Python with mpmath,
SymPy and cvxpy for numerics and certificates, Lean for one formalization, and
LaTeX built through a lint gate. This paper was drafted with Claude Code from
the project's records and logs and revised by the author. The agent
instructions, templates and tools are available in the accompanying process
repository.

\begin{bibdiv}
\begin{biblist}

\bib{alphaevolve}{article}{
  author={Novikov, A.},
  author={others},
  title={AlphaEvolve: a coding agent for scientific and algorithmic
      discovery},
  eprint={arXiv:2506.13131},
  date={2025},
}

\bib{alphageometry}{article}{
  author={Trinh, T. H.},
  author={Wu, Y.},
  author={Le, Q. V.},
  author={He, H.},
  author={Luong, T.},
  title={Solving olympiad geometry without human demonstrations},
  journal={Nature},
  volume={625},
  date={2024},
  pages={476\ndash 482},
  doi={10.1038/s41586-023-06747-5}, % chktex 8
}

\bib{avigad}{article}{
  author={Avigad, J.},
  title={Mathematics and the formal turn},
  journal={Bull. Amer. Math. Soc. (N.S.)},
  volume={61},
  date={2024},
  pages={225\ndash 240},
  doi={10.1090/bull/1832}, % chktex 8
}

\bib{bessis}{misc}{
  author={Bessis, D.},
  title={The fall of the theorem economy},
  date={2026},
  note={\url{https://davidbessis.substack.com/p/the-fall-of-the-theorem-economy}},
}

\bib{bubeck}{article}{
  author={Bubeck, S.},
  author={others},
  title={Early science acceleration experiments with GPT-5},
  eprint={arXiv:2511.16072},
  date={2025},
}

\bib{collins}{article}{
  author={Collins, K. M.},
  author={others},
  title={Evaluating language models for mathematics through
      interactions},
  eprint={arXiv:2306.01694},
  date={2023},
}

\bib{davies}{article}{
  author={Davies, A.},
  author={others},
  title={Advancing mathematics by guiding human intuition with AI},
  journal={Nature},
  volume={600},
  date={2021},
  pages={70\ndash 74},
  doi={10.1038/s41586-021-04086-x}, % chktex 8
}

\bib{erdos}{misc}{
  author={Bloom, T.},
  title={Erdős problems},
  note={\url{https://www.erdosproblems.com}},
}

\bib{erdos-wiki}{misc}{
  title={AI contributions to Erdős problems, and Disclaimers and
      caveats},
  note={Wiki of the Erdős problems repository,
      \url{https://github.com/teorth/erdosproblems/wiki}},
}

\bib{firstproof}{article}{
  author={Abouzaid, M.},
  author={Srivastava, N.},
  author={others},
  title={First Proof second batch},
  eprint={arXiv:2606.18119},
  date={2026},
}

\bib{frieder}{article}{
  author={Frieder, S.},
  author={others},
  title={Mathematical capabilities of ChatGPT},
  eprint={arXiv:2301.13867},
  date={2023},
}

\bib{frontiermath}{article}{
  author={Glazer, E.},
  author={others},
  title={FrontierMath: a benchmark for evaluating advanced mathematical
      reasoning in AI},
  eprint={arXiv:2411.04872},
  date={2024},
}

\bib{funsearch}{article}{
  author={Romera-Paredes, B.},
  author={others},
  title={Mathematical discoveries from program search with large
      language models},
  journal={Nature},
  volume={625},
  date={2024},
  pages={468\ndash 475},
  doi={10.1038/s41586-023-06924-6}, % chktex 8
}

\bib{georgiev}{article}{
  author={Georgiev, B.},
  author={Gómez-Serrano, J.},
  author={Tao, T.},
  author={Wagner, A. Z.},
  title={Mathematical exploration and discovery at scale},
  eprint={arXiv:2511.02864},
  date={2025},
}

\bib{gowers}{misc}{
  author={Gowers, W. T.},
  title={The two cultures of mathematics},
  note={\url{https://www.dpmms.cam.ac.uk/~wtg10/2cultures.pdf}},
}

\bib{hle}{article}{
  author={Phan, L.},
  author={others},
  title={Humanity's last exam},
  eprint={arXiv:2501.14249},
  date={2025},
}

\bib{htps}{article}{
  author={Lample, G.},
  author={others},
  title={HyperTree proof search for neural theorem proving},
  eprint={arXiv:2205.11491},
  date={2022},
}

\bib{leiden}{misc}{
  title={The Leiden Declaration on Artificial Intelligence and
      Mathematics},
  date={2026},
  note={\url{https://leidendeclaration.ai/},
      doi:10.5281/zenodo.20302944},
}

\bib{mathlib}{inproceedings}{
  author={{The mathlib Community}},
  title={The Lean mathematical library},
  booktitle={Proceedings of the 9th ACM SIGPLAN International Conference
      on Certified Programs and Proofs},
  date={2020},
  pages={367\ndash 381},
  doi={10.1145/3372885.3373824}, % chktex 8
}

\bib{polu}{article}{
  author={Polu, S.},
  author={Sutskever, I.},
  title={Generative language modeling for automated theorem proving},
  eprint={arXiv:2009.03393},
  date={2020},
}

\bib{quanta}{article}{
  author={Kakaes, K.},
  title={The AI revolution in math has arrived},
  journal={Quanta Magazine},
  date={2026},
  note={\url{https://www.quantamagazine.org/the-ai-revolution-in-math-has-arrived-20260413/}},
}

\bib{tao-age}{article}{
  author={Tao, T.},
  title={Mathematics in the age of AI},
  eprint={arXiv:2608.16753},
  date={2026},
}

\bib{tao-cleverness}{misc}{
  author={Tao, T.},
  title={Post on artificial general cleverness},
  date={2025},
  note={Mathstodon, \url{https://mathstodon.xyz/@tao/115722360006034040}},
}

\bib{tao-good}{article}{
  author={Tao, T.},
  title={What is good mathematics?},
  journal={Bull. Amer. Math. Soc. (N.S.)},
  volume={44},
  date={2007},
  pages={623\ndash 634},
  eprint={arXiv:math/0702396},
}

\bib{tao-lowhanging}{misc}{
  author={Tao, T.},
  title={Post on AI solutions to Erdős problems},
  date={2025},
  note={Mathstodon, \url{https://mathstodon.xyz/@tao/115788262274999408}},
}

\bib{thurston}{article}{
  author={Thurston, W. P.},
  title={On proof and progress in mathematics},
  journal={Bull. Amer. Math. Soc. (N.S.)},
  volume={30},
  date={1994},
  pages={161\ndash 177},
  eprint={arXiv:math/9404236},
}

\bib{wagner}{article}{
  author={Wagner, A. Z.},
  title={Constructions in combinatorics via neural networks},
  eprint={arXiv:2104.14516},
  date={2021},
}

\end{biblist}
\end{bibdiv}

\end{document}
